\chapter*{Preface}
\label{frontmatter:preface}
\addcontentsline{toc}{chapter}{Preface}

A stone, a rope, a fire, a lever, a stream, a wall, a seed, a signal, a tool. Engineering begins there. Everything that follows in this book is the disciplined continuation of those primitive acts. The reader who works with us through the twelve volumes learns to lift, join, heat, flow, sense, count, grow, control, remember, and repair, with the precision and the responsibility the artificial world demands.

We address this book to one declared audience. It is the serious adult learner, with the imagined reader being a capable twenty-five-year-old who wants to become an engineer in mind and practice. We speak to that reader directly, while allowing older professionals, self-taught builders, and students to overhear without feeling excluded. We do not assume the reader is already an engineer. We assume the reader is willing to work, willing to be wrong, and willing to redo the calculation.

We have a thesis. It is the sentence that governs the twelve volumes:

\begin{quote}\itshape
Engineering is the discipline by which measured reality becomes reliable intervention under constraint, failure, scale, and responsibility.
\end{quote}

That is what we mean by \emph{Engineering} with a capital E. Engineering is not applied science. Engineering is disciplined intervention. Science asks what is true. Engineering asks what can be made to work, for whom, under what constraints, for how long, and at what cost if it fails.

This book is not an encyclopedia. An encyclopedia tries to include every fact. We try to include every major class of problem, and to teach the reader to reduce new problems to known structures. The reader who solves several thousand carefully chosen archetypal problems can attack millions of specific ones. We give grammar; we do not give vocabulary.

The book is sequential, not modular. We read it front to back at least once. The twelve volumes, taken in order, are the twelve stages of a single ascent: from quantity (what we can measure) through form (what we can model), force, energy, matter, life, information, machines, systems, failure, design, and at the end civilization itself. Each volume builds on every previous volume. The reader who skips Volume III on Force will read Volume VIII on Machines as a stranger.

We cannot make a Professional Engineer of the reader. Professional licensure depends on accredited education, supervised experience, and examinations, and we do not pretend to substitute for those institutions. We can give the reader the substantive ability that licensure attests to. Credential and capability are different objects, and a serious reader will eventually want both.

What follows is structured as a story. Reality is measurable. Measurement becomes model. Model becomes mechanism. Mechanism becomes system. System breaks. Breakage becomes discipline. Discipline becomes civilization. The story has a beginning. It does not have a final chapter, because the artificial world is not finished and the reader is the one who continues it.

Read carefully. Solve the exercises. Build the projects. Document the failures. Trust no model that has not been humiliated by reality.

That last instruction is half of engineering.
